\scanpage{39}\mspage{45}
\deleted{Un prédérivateur est dit}
% slanted
\margin{Inutile comme axiome, \unsure{va|car} \unsure{devenir|devient} un théorème (résulte des autres axiomes}
\uline{Der 4}$_0$ (g) : Soit
\diagram{\begin{tikzcd}[ampersand replacement=\&]
 U \arrow[r,hook,"j"'] \& X \& Y \arrow[l,hook',"i"]
\end{tikzcd}}
\inserted{dans Diag}, \unsure{imm|immersions} ouverte et fermée complémentaires l'une de l'autre. Alors le morphisme d'adjonction dans $\mathbb D(Y)$
\[
 i^*i_*(F_Y)\longrightarrow F_Y
\]
est \deleted{toujours \unsure{plein}} \inserted{un iso}, donc \deleted{\unsure{$i_*$ est}}
\[
 i_*:\mathbb D(Y)\lhook\joinrel\longrightarrow\mathbb D(X)
\]
est pleinement fidèle, et $i^*$\inserted{$:\mathbb D(X)\to\unsure{\mathbb D(Y)|\mathbb D Y}$} est un foncteur de localisation. \deleted{Plus précisément, pour \inserted{$u:F\to G$} dans \unsure{$I$|$J$}, \deleted{\unsure{si $u$}} \unsure{$u$} est un iso ssi $\unsure{i|i^*}$ $i^*(u):i^*(F)\to i^*(G)$ \unsure{est un iso}} L'image essentielle \inserted{de $i_*$} est formée des objets $F$ de $\mathbb D(X)$ tels que $j^*(F)$ soit l'objet final de $\mathbb D(U)$.

\uline{Der \unsure{4|5}}$_0$ (d) Avec les notations précédentes pour $i,j$, le morphisme d'adjonction dans $\mathbb D(\unsure{u|U})$
\[
 F_U\to j_{\deleted{!}}^*j_{\unsure{!}}^{\deleted{\unsure{*}}}F_{\unsure{u|U}}\ \deleted{\unsure{\to F_U}}
\]
est un iso -- $j_{\unsure{*|!}}$ est pl. fid, $j^*$ un foncteur localisation. L'image essentielle de $j_!$ est la sous-catégorie pleine de $\mathbb D(X)$ formée des $F$ \deleted{\unsure{de}} tels que $i^*(F)$ soit \deleted{\unsure{l'}}\unsure{obj} initial \deleted{\unsure{de}} de $\mathbb D(Y)$.

Il y a les deux carrés de la page 8 et 9, mais pour interpréter \unsure{les|le} \unsure{caractères|caractère|condition} cartésien ou cocartésien c'est prématuré ici.
% archive sheet number: 38
